\documentclass[12pt,a4paper]{article}
\usepackage[margin=1in]{geometry}
\usepackage{amsmath}
\usepackage{graphicx}
\usepackage{hyperref}
\usepackage{booktabs}
\usepackage{array}
\usepackage{enumitem}
\usepackage{parskip}
\usepackage{titlesec}
\usepackage{fancyhdr}
\usepackage{xcolor}
\usepackage{listings}
\usepackage{caption}
\usepackage{float}

\hypersetup{
    colorlinks=true,
    linkcolor=blue,
    urlcolor=blue,
    citecolor=blue
}

\lstset{
    basicstyle=\ttfamily\small,
    breaklines=true,
    frame=single,
    backgroundcolor=\color{gray!10},
    keywordstyle=\color{blue}\bfseries,
    commentstyle=\color{green!50!black},
    stringstyle=\color{red!70!black}
}

\pagestyle{fancy}
\fancyhf{}
\rhead{Restaurant Management System}
\lhead{CSC 4710 --- Database Systems}
\cfoot{\thepage}

\title{
    \vspace{-1cm}
    \Large\textbf{Restaurant Management System} \\[6pt]
    \large CSC 4710 --- Database Systems \\[4pt]
    \normalsize Georgia State University
}
\author{
    Raphael Omorose \and Ahmed Cisse \and Jonathan Bell
}
\date{April 2026}

\begin{document}

\maketitle
\thispagestyle{fancy}

\tableofcontents
\newpage

% -----------------------------------------------------------------------
\section{Introduction}
% -----------------------------------------------------------------------

\subsection{Application Overview}

The Restaurant Management System (RMS) is a full-stack web application that
digitalises the daily operations of a restaurant. It supports two distinct user
roles: \textbf{customers}, who browse the menu, build a cart, and place orders;
and \textbf{administrators}, who manage ingredients, menu items, inventory stock
levels, order fulfilment, and business analytics --- all from a single
browser-based interface with no specialist software required.

\subsection{Motivation}

We chose a restaurant management system because it exercises every major database
concept covered in the course in a context that is immediately intuitive to anyone
who has visited a restaurant.

\begin{itemize}
    \item \textbf{Multi-table relational design.} Ingredients, menu items, orders,
          users, and inventory are naturally distinct entities with meaningful
          relationships, requiring thoughtful schema design.

    \item \textbf{Integrity constraints.} Price must be positive, ingredient names
          cannot be blank, and an ingredient referenced by a menu item must not be
          deleted. These constraints are enforced at both the database level (NOT
          NULL, FK cascades) and the application layer (bean validation,
          ConflictException).

    \item \textbf{Transactions.} Placing an order must atomically create
          order-item rows, deduct ingredient inventory, and record a price
          snapshot. Partial success is never acceptable and is prevented with
          \texttt{@Transactional}.

    \item \textbf{Aggregate queries.} Revenue trends, top-selling items, order
          volume, and inventory forecasting are meaningful business questions
          that require \texttt{GROUP BY}, \texttt{SUM}, and date arithmetic ---
          direct applications of the SQL taught in this course.

    \item \textbf{Normalisation.} The recipe junction table
          (\texttt{menu\_item\_ingredients}) cleanly separates the many-to-many
          relationship between menu items and ingredients. The price snapshot on
          \texttt{order\_items} avoids update anomalies from future price changes.
\end{itemize}

\subsection{Key Components}

\begin{itemize}
    \item \textbf{Spring Boot 4 REST API} (Java 21) --- nine controllers
          exposing approximately 40 endpoints.
    \item \textbf{MySQL 8 database} --- eight normalised tables managed via
          Spring Data JPA / Hibernate.
    \item \textbf{React 19 / TypeScript SPA} --- five pages, eleven reusable
          admin-tab components, and a Recharts analytics dashboard.
    \item \textbf{Axios HTTP client} --- typed API layer split into five modules.
    \item \textbf{Cloudinary} --- image uploads compressed server-side and stored
          in Cloudinary; the API returns a \texttt{secure\_url} saved in the
          database.
    \item \textbf{Spring Security} --- full security filter chain with CSRF
          disabled, stateless sessions, and CORS delegated to \texttt{WebConfig}.
          BCrypt password hashing; no plaintext passwords are ever stored or
          logged. Role enforcement is applied at the service layer and via
          frontend route guards.
\end{itemize}

% -----------------------------------------------------------------------
\section{Database Design}
% -----------------------------------------------------------------------

\subsection{Entity--Relationship Model}

The ER model contains eight entities and the following relationships.

\begin{itemize}
    \item \textbf{User} places zero or more \textbf{Order}s \hfill(1:N)
    \item \textbf{Order} contains one or more \textbf{OrderItem}s \hfill(1:N)
    \item Each \textbf{OrderItem} references exactly one \textbf{MenuItem} \hfill(N:1)
    \item \textbf{MenuItem} is composed of zero or more \textbf{Ingredient}s via
          the associative entity \textbf{MenuItemIngredient} \hfill(M:N $\to$ two 1:N)
    \item Each \textbf{Ingredient} has exactly one \textbf{IngredientInventory}
          record (1:1 shared-key; \texttt{ingredient\_inventory.ingredient\_id}
          is both PK and FK)
    \item \textbf{StockChangeLog} references the acting admin \textbf{User} and
          the affected \textbf{Ingredient} \hfill(N:1 each)
\end{itemize}

\subsection{Relational Schema}

Table~\ref{tab:schema} summarises the eight tables, their primary keys, and
their foreign key dependencies.

\begin{table}[H]
\centering
\caption{Database schema overview}
\label{tab:schema}
\renewcommand{\arraystretch}{1.3}
\begin{tabular}{@{}p{4cm}p{3.2cm}p{3.5cm}p{3.5cm}@{}}
\toprule
\textbf{Table} & \textbf{Primary Key} & \textbf{Foreign Keys} & \textbf{Notable Constraints} \\
\midrule
\texttt{users}
    & \texttt{user\_id}
    & ---
    & \texttt{username} UNIQUE; \texttt{role} ENUM(\texttt{USER},\texttt{ADMIN}) \\

\texttt{ingredients}
    & \texttt{ingredient\_id}
    & ---
    & \texttt{name}, \texttt{unit} NOT NULL; \texttt{image\_url} pattern-validated \\

\texttt{ingredient\_inventory}
    & \texttt{ingredient\_id}
    & \texttt{ingredient\_id} $\to$ \texttt{ingredients}
    & 1:1 shared PK/FK; \texttt{quantity\_on\_hand} $\geq 0$ \\

\texttt{menu\_items}
    & \texttt{menu\_item\_id}
    & ---
    & \texttt{price} $\geq 0.01$; \texttt{is\_active} DEFAULT TRUE;
      \texttt{category} ENUM (BURGERS, SANDWICHES, FRIES, SIDES, BREAKFAST,
      SOUPS, SWEETS, DRINKS) \\

\texttt{menu\_item\_ingredients}
    & (\texttt{menu\_item\_id}, \texttt{ingredient\_id})
    & both cols are FKs
    & \texttt{quantity\_required} $> 0$ \\

\texttt{orders}
    & \texttt{order\_id}
    & \texttt{user\_id} $\to$ \texttt{users}
    & \texttt{status} $\in$ \{PLACED, READY, COMPLETE, CANCELLED\} \\

\texttt{order\_items}
    & \texttt{order\_item\_id}
    & \texttt{order\_id}, \texttt{menu\_item\_id}
    & \texttt{unit\_price\_at\_time} immutable snapshot \\

\texttt{stock\_change\_log}
    & \texttt{stock\_change\_id}
    & \texttt{admin\_user\_id}, \texttt{ingredient\_id}
    & Append-only audit trail \\
\bottomrule
\end{tabular}
\end{table}

\subsection{Functional Dependencies and Normalisation}

All eight tables satisfy \textbf{Boyce--Codd Normal Form (BCNF)}.

\begin{itemize}
    \item Every table with a single-column primary key has all non-key attributes
          functionally dependent \emph{solely} on that key --- no partial or
          transitive dependencies exist.

    \item \texttt{menu\_item\_ingredients} uses a composite primary key
          (\texttt{menu\_item\_id}, \texttt{ingredient\_id}). Its only non-key
          attribute, \texttt{quantity\_required}, depends on the \emph{full}
          composite key, satisfying BCNF.

    \item \texttt{ingredient\_inventory} shares its primary key with
          \texttt{ingredients} (1:1 shared-key pattern). The dependency
          \texttt{ingredient\_id} $\to$ \{\texttt{quantity\_on\_hand},
          \texttt{last\_updated\_at}\} is trivially BCNF.

    \item \texttt{order\_items} stores \texttt{unit\_price\_at\_time} as a
          deliberate event snapshot. Although
          \texttt{menu\_item\_id} $\to$ \texttt{price} holds in
          \texttt{menu\_items}, recording the price \emph{at the time of the
          order} is a fact about the order event itself, not a derived value.
          Historical orders must not be affected by future price changes.
          This is a sound design decision, not a normalisation violation.
\end{itemize}

\subsection{Additional Constraints}

\begin{itemize}
    \item \textbf{Application-level referential protection.} Deleting an
          ingredient that is used by any menu item raises HTTP 409 Conflict
          (\texttt{ConflictException}) before any database statement is executed.

    \item \textbf{Bean validation.} \texttt{@NotBlank} on ingredient and
          menu-item names, \texttt{@DecimalMin("0.01")} on price, and
          \texttt{@Pattern} on image URLs are enforced via \texttt{@Valid} on
          controller request bodies before any write reaches the database.

    \item \textbf{Password security.} BCrypt via Spring Security Crypto. The
          raw password is write-only (\texttt{@JsonProperty(access =
          WRITE\_ONLY)}) and never appears in API responses or logs.

    \item \textbf{Audit trail.} \texttt{stock\_change\_log} is append-only ---
          no update or delete endpoints are exposed, preserving a full history
          of every inventory change. Entries are written both when an admin
          manually adjusts stock and when \texttt{placeOrder} deducts
          ingredients, recording the ordering user as the acting party.

    \item \textbf{Performance indexes.} JPA \texttt{@Index} annotations on
          \texttt{orders(placed\_at)}, \texttt{orders(status)},
          \texttt{order\_items(menu\_item\_id)}, and
          \texttt{order\_items(order\_id)} are applied automatically by
          Hibernate on startup (\texttt{ddl-auto=update}). These indexes
          directly benefit the analytics native SQL queries which filter and
          group on \texttt{placed\_at} and join on \texttt{menu\_item\_id}.
\end{itemize}

% -----------------------------------------------------------------------
\section{Functionality Details}
% -----------------------------------------------------------------------

\subsection{Basic Functions}

\begin{description}[style=nextline]
    \item[User registration and authentication]
    \texttt{POST /api/users/register} hashes the password with BCrypt and
    persists the user. \texttt{POST /api/users/login} verifies the hash and
    returns the user object (including role) for client-side session storage.
    The frontend stores the response in \texttt{localStorage} and guards
    routes accordingly; non-admin users are redirected away from
    \texttt{/admin}.

    \item[Menu browsing and availability]
    \texttt{GET /api/menu-items/availability} joins menu items with their
    ingredient requirements and current inventory levels to return a real-time
    availability flag for each dish. The customer order page uses this to
    visually indicate out-of-stock items before the customer tries to order.

    \item[Cart and order placement]
    Customers add items to a client-side cart (React state). On checkout,
    \texttt{POST /api/orders/place} runs a \texttt{@Transactional} method that:
    (1) validates the cart is non-empty;
    (2) validates all menu items exist and are active;
    (3) creates the \texttt{Order} and \texttt{OrderItem} rows with a price
    snapshot; and
    (4) deducts ingredient inventory for every ordered item across all
    ingredient rows.

    \item[Order history]
    \texttt{GET /api/orders/user/\{userId\}} returns the customer's full order
    history with items eagerly fetched via a JPQL \texttt{JOIN FETCH} query,
    ordered by \texttt{created\_at DESC} to show the most recent orders first.

    \item[Image upload]
    \texttt{POST /api/images/upload} accepts a \texttt{multipart/form-data}
    file, compresses it server-side via the Cloudinary Java SDK
    (quality auto, format auto, 800\,px width limit), and returns the
    resulting \texttt{secure\_url}. Ingredient and menu-item forms in the admin
    panel use this endpoint to replace the plain URL text input with a
    file-picker button and live thumbnail preview.

    \item[Ingredient and menu item CRUD]
    Full create, read, update, and delete for ingredients (\texttt{/api/ingredients})
    and menu items (\texttt{/api/menu-items}) via standard REST endpoints.
    Menu items carry a \texttt{category} field (one of eight diner categories)
    displayed and filtered in the admin panel.
    Ingredient deletes are protected against referential use.

    \item[Inventory management]
    Admins view and update \texttt{quantity\_on\_hand} for each ingredient via
    \texttt{IngredientInventory} endpoints. Every manual adjustment can be
    recorded to \texttt{stock\_change\_log} for audit purposes.
\end{description}

\subsection{Advanced Functions}

\begin{description}[style=nextline]
    \item[Order status lifecycle]
    Orders follow a process-driven lifecycle rather than a manually settable
    status field. When \texttt{POST /api/orders/place} is called, the order
    enters the \texttt{PLACED} state. A grace window of 5--10 seconds
    (proportional to item count) allows the customer to cancel via
    \texttt{POST /api/orders/\{id\}/cancel}. After the window expires, the
    next read of that order triggers an in-process transition to
    \texttt{COMPLETE}. The admin Orders tab auto-refreshes every two seconds,
    shows a live countdown for \texttt{PLACED} orders, and supports status
    filtering via \texttt{GET /api/orders?status=\{status\}}, returning
    HTTP~400 for unrecognised values.

    \item[Filtered order queries]
    \texttt{GET /api/orders?status=PLACED} uses a Spring Data derived query
    (\texttt{findByStatus}) to return only orders in a given state, enabling
    the admin to focus on pending work without client-side filtering.

    \item[Analytics dashboard]
    Four native SQL queries (via JPA \texttt{EntityManager}) power the
    analytics tab.
    \begin{itemize}
        \item \textbf{Revenue over time} --- daily \texttt{SUM(quantity
              $\times$ unit\_price\_at\_time)} grouped by
              \texttt{DATE(placed\_at)}, rendered as a Recharts
              \texttt{LineChart}.
        \item \textbf{Top-selling items} --- total units and revenue per menu
              item in a date range, rendered as a \texttt{BarChart}.
        \item \textbf{Order volume} --- daily order counts in a date range,
              rendered as a \texttt{LineChart}.
        \item \textbf{Inventory forecast} --- joins ingredient demand (from
              recent orders) against current stock to estimate remaining supply,
              rendered as a \texttt{BarChart}.
    \end{itemize}

    \item[Checkout and order summary flow]
    The customer checkout is a two-step flow. The cart panel shows a running
    total and a ``Checkout'' button that navigates to a dedicated Order Summary
    page. The summary page lists every cart item with its thumbnail image,
    quantity, and line total alongside a full pricing breakdown (subtotal,
    8\% tax, service fee, and total). Confirming calls
    \texttt{POST /api/orders/place}; on success a CSS keyframe animation draws
    an orange SVG checkmark (stroke-dashoffset from 40 to 0) inside a
    ring-draw circle, accompanied by a spring-bounce entrance. The customer
    can then navigate to their order history or return to the menu.
\end{description}

% -----------------------------------------------------------------------
\section{Implementation Details}
% -----------------------------------------------------------------------

\subsection{Technology Stack}

\begin{table}[H]
\centering
\caption{Technology choices}
\label{tab:tech}
\renewcommand{\arraystretch}{1.2}
\begin{tabular}{@{}lll@{}}
\toprule
\textbf{Layer} & \textbf{Technology} & \textbf{Version} \\
\midrule
Database            & MySQL                          & 8.0 \\
ORM / persistence   & Spring Data JPA / Hibernate    & via Spring Boot 4 \\
Backend framework   & Spring Boot                    & 4.0.3 \\
Backend language    & Java                           & 21 \\
Build tool          & Apache Maven                   & 3.x \\
Frontend framework  & React                          & 19.2 \\
Frontend language   & TypeScript                     & 5.9 \\
Frontend build tool & Vite                           & 7.3 \\
HTTP client         & Axios                          & 1.15 \\
Charting            & Recharts                       & 3.8 \\
Client-side routing & React Router                   & 7.13 \\
Security framework  & Spring Security                & via Spring Boot 4 \\
Password hashing    & BCrypt (Spring Security)       & via Spring Boot 4 \\
\bottomrule
\end{tabular}
\end{table}

\subsection{Backend Architecture}

The backend follows a standard four-layer architecture.

\begin{enumerate}
    \item \textbf{Entity layer.} JPA-annotated POJOs with Lombok for boilerplate
          reduction. Jackson annotations
          (\texttt{@JsonManagedReference}/\texttt{@JsonBackReference}) prevent
          circular serialisation on bidirectional relationships
          (\texttt{Order $\leftrightarrow$ OrderItem},
          \texttt{MenuItem $\leftrightarrow$ MenuItemIngredient}).

    \item \textbf{Repository layer.} Spring Data JPA interfaces. Custom JPQL
          queries handle \texttt{JOIN FETCH} (eliminating N+1 selects) and
          derived queries for filtered lookups. Native SQL is used only where
          JPQL cannot express the required aggregation.

    \item \textbf{Service layer.} All business logic lives here: transaction
          boundaries (\texttt{@Transactional}), input validation, inventory
          deduction, status transition validation, and the referential delete
          guard.

    \item \textbf{Controller layer.} Thin REST controllers that map HTTP verbs
          to service methods. Bean validation is triggered via \texttt{@Valid}
          on request bodies. A \texttt{@RestControllerAdvice} global handler
          maps \texttt{ResourceNotFoundException} to 404 and
          \texttt{ConflictException} to 409 with a consistent JSON error body.
\end{enumerate}

CORS is configured in \texttt{WebConfig} to allow the Vite dev server
(\texttt{localhost:5173}) to call the Spring Boot server
(\texttt{localhost:8080}) using all standard HTTP methods.
\texttt{SecurityConfig} delegates CORS decisions to \texttt{WebConfig} via
\texttt{Customizer.withDefaults()}, ensuring the Spring Security filter chain
does not shadow the MVC-level CORS mappings.

\subsection{Frontend Architecture}

The frontend is a single-page application using React Router v7 with five
routes.

\begin{itemize}
    \item \texttt{/} --- Landing page with feature overview and navigation to
          both login flows.
    \item \texttt{/login} --- Combined sign-in / sign-up form. Redirects to
          \texttt{/order} for customers and \texttt{/admin} for admins based
          on the role returned by the login endpoint.
    \item \texttt{/admin-login} --- Dedicated admin portal that rejects
          non-admin credentials at the client level.
    \item \texttt{/order} --- Customer page with a Menu tab (live availability,
          CartPanel checkout) and a My Orders tab (order history).
    \item \texttt{/admin} --- Admin page with six tabs: Ingredients, Menu Items,
          Inventory, Stock Log, Orders, and Analytics.
\end{itemize}

The API layer (\texttt{src/api/}) is split into six typed Axios modules.

\begin{itemize}
    \item \texttt{axios.ts} --- base instance pointed at
          \texttt{http://localhost:8080}.
    \item \texttt{adminApi.ts} --- CRUD for ingredients, menu items, inventory,
          and stock log; also exposes \texttt{uploadImage()} which posts a
          \texttt{File} to \texttt{/api/images/upload} and returns the
          Cloudinary URL.
    \item \texttt{menuApi.ts} --- public menu and availability reads.
    \item \texttt{orderApi.ts} --- admin order management (list all, filter by
          status, cancel); also contains the shared \texttt{Order} type and
          cancellation-window calculation used by both admin and customer views.
    \item \texttt{analyticsApi.ts} --- four analytics endpoints with date-range
          parameters; \texttt{getAllAnalytics()} fetches all four in parallel
          via \texttt{Promise.all}.
    \item \texttt{customerOrderApi.ts} --- customer order placement
          (\texttt{POST /api/orders/place}) and per-user history.
\end{itemize}

\subsection{Frontend--Backend Interaction}

The frontend communicates with the backend exclusively over HTTP/JSON on
\texttt{localhost:8080/api/**}. Authentication is stateless from the server's
perspective: the login endpoint returns a user object that the browser stores in
\texttt{localStorage}. Route guards in each page component read this value on
mount and redirect unauthenticated or unauthorised users. This approach was
chosen for simplicity in an academic context; a production system would use
short-lived JWTs and server-side session validation on every protected endpoint.

\subsection{Source Code Repository}

The complete source code is publicly available at:

\begin{center}
\url{https://github.com/OfficialEseosa/Restaurant-Management-System}
\end{center}

% -----------------------------------------------------------------------
\section{Group Member Evaluation}
% -----------------------------------------------------------------------

\begin{table}[H]
\centering
\caption{Group member contributions and scores}
\label{tab:team}
\renewcommand{\arraystretch}{1.4}
\begin{tabular}{@{}p{3.2cm}p{9cm}c@{}}
\toprule
\textbf{Member} & \textbf{Contributions} & \textbf{Score (/10)} \\
\midrule
Raphael Omorose
    & Backend lead. Designed the full database schema. Built all service and
      repository layers including analytics native SQL queries, menu availability
      endpoint, atomic order placement with inventory deduction, order status
      management, and CORS configuration.
    & 10 \\

Ahmed Cisse
    & Backend validation and frontend integration. Added Bean Validation
      annotations to entities, \texttt{@Valid} controller enforcement,
      \texttt{ConflictException} with global error handler, ingredient delete
      guard, and contributed to the admin panel frontend components.
    & 9 \\

Jonathan Bell
    & Frontend integration. Built the complete React/TypeScript SPA: all five
      pages, eleven admin-tab components (Ingredients, Menu Items, Inventory,
      Stock Log, Orders, Analytics), the CartPanel checkout flow, OrderHistory,
      and all five typed Axios API modules.
    & 10 \\
\bottomrule
\end{tabular}
\end{table}

% -----------------------------------------------------------------------
\section{Experiences}
% -----------------------------------------------------------------------

\subsection{Lessons Learned}

\textbf{Normalisation has real practical consequences.}
Storing \texttt{unit\_price\_at\_time} as a snapshot on \texttt{order\_items}
instead of joining back to \texttt{menu\_items.price} was initially debated as a
normalisation violation. Working through the functional dependency argument
clarified that it is not: the fact being recorded is ``the price charged at the
time of this order'', which depends on the order event, not on the menu item
alone. Without the snapshot, any price update would retroactively change
historical order totals --- a serious data integrity problem.

\textbf{N+1 queries are silent until they hurt.}
Because Hibernate's default behaviour fetches associated collections lazily
(or eagerly via a separate \texttt{SELECT} per row), a \texttt{findAll()} over
23 ingredients was issuing 24 round-trips to a remote database --- one for the
ingredient list and one per row for the joined \texttt{ingredient\_inventory}
record. At roughly 150\,ms per round-trip this produced a 3--5 second response
time. Overriding \texttt{findAll()} on every affected repository with
\texttt{@EntityGraph} collapses all associations into a single \texttt{LEFT
JOIN}, cutting each list endpoint to a single database call regardless of row
count.

\textbf{Transactionality must be designed, not assumed.}
The \texttt{placeOrder} method creates an order, multiple order items, and then
deducts inventory across potentially many ingredient rows. Wrapping the entire
method in \texttt{@Transactional} ensures that a failure at any point (e.g., a
missing inventory record) rolls back all changes atomically, leaving the database
in a consistent state.

\textbf{JPA \texttt{@MapsId} requires \texttt{persist}, not \texttt{merge}.}
The \texttt{POST /api/ingredients} endpoint returned HTTP~500 with a Hibernate
\texttt{AssertionFailure: null identifier} on every call. The root cause was
subtle: \texttt{ensureInventoryRow} explicitly set \texttt{inventory.setIngredientId(id)}
before calling \texttt{repository.save()}. Spring Data JPA interprets a
non-null \texttt{@Id} field as evidence of a detached entity and calls
\texttt{em.merge()} instead of \texttt{em.persist()}. During merge, Hibernate
traverses the \texttt{@MapsId} association and fails to derive the shared
primary key from a proxy obtained in a different transaction. The fix was
twofold: remove the explicit \texttt{setIngredientId()} call (letting
\texttt{@MapsId} derive the PK from the \texttt{ingredient} reference), and
annotate \texttt{IngredientService.save()} with \texttt{@Transactional} to
keep all operations in a single \texttt{EntityManager} session.

\subsection{Hard Problems Solved}

\begin{itemize}
    \item \textbf{Circular JSON serialisation.} Bidirectional JPA relationships
          (\texttt{Order $\leftrightarrow$ OrderItem},
          \texttt{MenuItem $\leftrightarrow$ MenuItemIngredient}) cause infinite
          loops during Jackson serialisation. Resolved with
          \texttt{@JsonManagedReference} / \texttt{@JsonBackReference} pairs so
          only one side of each relationship is serialised.

    \item \textbf{Composite primary key mapping.} The recipe junction table uses
          a composite PK. Mapping this correctly required an
          \texttt{@Embeddable} ID class (\texttt{MenuItemIngredientId}) and
          \texttt{@EmbeddedId} on the entity, with \texttt{@MapsId} on each FK
          field.

    \item \textbf{Analytics with native SQL.} The aggregate queries for revenue,
          top items, and inventory forecast could not be expressed in JPQL.
          Using \texttt{EntityManager.createNativeQuery} with
          \texttt{@PersistenceContext} injection provided full SQL
          expressiveness while remaining within the Spring/JPA ecosystem.

    \item \textbf{Spring Security shadowing MVC CORS.} Adding
          \texttt{spring-boot-starter-security} inserted a filter chain that
          ran before the MVC dispatcher. Without explicit configuration,
          Spring Security's \texttt{CorsFilter} rejected cross-origin preflight
          requests before they reached \texttt{WebConfig}'s CORS mappings.
          Resolved by calling \texttt{.cors(Customizer.withDefaults())} in
          \texttt{SecurityConfig}, which delegates CORS decisions back to the
          MVC-level configuration rather than blocking at the security filter
          layer.

    \item \textbf{Cloudinary multipart upload with server-side compression.}
          Storing images required a dedicated \texttt{POST /api/images/upload}
          endpoint that receives a \texttt{MultipartFile}, uploads the raw bytes
          to the Cloudinary API with an eager transformation
          (\texttt{q\_auto,f\_auto,w\_800,c\_limit}), and returns the
          \texttt{secure\_url}. The frontend file-picker sends a
          \texttt{FormData} request and immediately renders a thumbnail from
          the returned URL before the user saves the form.

    \item \textbf{Transient entity reference on composite-key save.}
          The \texttt{POST /api/menu-item-ingredients} endpoint consistently
          returned HTTP~400. Jackson deserialises the request body into a
          \texttt{MenuItemIngredient} whose \texttt{MenuItem} and
          \texttt{Ingredient} fields are transient stubs --- only their IDs are
          set. With no cascade configuration, Hibernate attempted to persist
          these unmanaged objects and threw
          \texttt{TransientPropertyValueException}. The fix was to replace the
          Jackson-created stubs with JPA proxy references from
          \texttt{getReferenceById()} and to explicitly construct the
          \texttt{@EmbeddedId} in the service layer before calling
          \texttt{save()}.
\end{itemize}

\subsection{Future Extensions}

\begin{itemize}
    \item \textbf{JWT authentication.} The Spring Security filter chain is
          already in place with a permissive configuration. The natural next
          step is to add a JWT filter: issue a signed token on login, require
          \texttt{Authorization: Bearer} on protected endpoints, and replace
          \texttt{localStorage} role storage with server-enforced claims.
    \item \textbf{Real-time order notifications.} Use WebSockets or
          Server-Sent Events so the kitchen view updates when new orders arrive,
          eliminating the need for manual page refreshes.
    \item \textbf{Multi-location support.} Extend the schema with a
          \texttt{locations} table and add \texttt{location\_id} foreign keys
          to inventory and orders to support restaurant chains.
    \item \textbf{Payment integration.} Integrate a payment gateway (e.g.,
          Stripe); add a \texttt{payments} table with idempotency keys to
          handle retry safety and prevent double-charges.
\end{itemize}

% -----------------------------------------------------------------------
\section{References}
% -----------------------------------------------------------------------

\begin{enumerate}
    \item Spring Boot Documentation. \textit{Spring Boot Reference Guide 4.x}.
          \url{https://docs.spring.io/spring-boot/docs/current/reference/html/}

    \item Spring Data JPA Documentation. \textit{Spring Data JPA Reference
          Documentation}.
          \url{https://docs.spring.io/spring-data/jpa/docs/current/reference/html/}

    \item React Documentation. \textit{React 19 --- Learn React}.
          \url{https://react.dev}

    \item Recharts. \textit{Recharts --- Redefined chart library built with React
          and D3}. \url{https://recharts.org}

    \item React Router. \textit{React Router v7 Documentation}.
          \url{https://reactrouter.com}

    \item Axios. \textit{Axios HTTP Client Documentation}.
          \url{https://axios-http.com}

    \item MySQL 8.0 Reference Manual. Oracle Corporation.
          \url{https://dev.mysql.com/doc/refman/8.0/en/}

    \item Elmasri, R.\ \& Navathe, S.\ B\@. \textit{Fundamentals of Database
          Systems}, 7th ed. Pearson, 2016.

    \item OWASP\@. \textit{OWASP Top Ten --- A07:2021 Identification and
          Authentication Failures}. \url{https://owasp.org/Top10/}

    \item Vite. \textit{Vite --- Next Generation Frontend Tooling}.
          \url{https://vite.dev}
\end{enumerate}

\end{document}
